\documentclass{article}
\usepackage{amsmath,hyperref}
\begin{document}
\section*{A relay sized for a scheduled cut need not cover an unexpected loss}
A single standby with the original range of 1.65 can cover each of two abrupt-loss identities separately in the first state of the published three-dimensional formation. It cannot cover both possibilities from one common pre-fault position, even with collision clearance removed. With the initial configuration unchanged, two appropriately placed ordinary-range standbys can cover them at the cut, or the one relay needs range about 1.791923. This is a resource lower bound, not delay in diagnosing which member failed.

\paragraph{Enlarged fault support.}
Use the same pinned \href{https://github.com/KTH-DHSG/constrained_omrs/tree/cb87f91aaaa4e652de043704ce4dfb989e10bcde}{author point-mass simulation} and \href{https://github.com/eternal-november/eos-research-programmes/pull/471}{standby check} on the nominal three-dimensional trajectory. Here the admissible cuts also include earlier losses. For each hypothesis only the selected vehicle's communication is removed; its physical body remains for clearance. Other poses and their possible geometric links are those of the recorded nominal state at that cut. Before the fault no new information distinguishes the two initial loss identities. No trajectory after either fault is taken from the source experiment. Original authorized formation links may be fewer than the full distance graph used here, making the impossibility for that full graph the stronger negative statement. Positive statements assume the advertised original and relay links are available and activated. Distance and clearance units are the point-mass model's units.

At time zero the loss of robot 1 leaves components $\{0\}$ and $\{2,3,4\}$, and the loss of robot 3 leaves $\{0,1,2\}$ and $\{4\}$. A common relay must therefore reach original robots 0 and 4 at their time-zero positions. These robots lie in different components in each history, so their two links are also sufficient. The radius is bounded below by half their separation and attained by their midpoint before testing clearance. The midpoint in this case clears all original bodies. Independently placing a relay at the midpoint of 0 and 2 for the first case, or 2 and 4 for the second, works at 1.65. Keeping both simultaneously is allowed by the pairwise clearance model; their mutual distance and distance to all original bodies exceed 0.45. Neither assignment moves in response to a late observation. The alternative of rerouting original robots before loss changes the pre-fault configurations and is not ruled out by this lower bound. Likewise a mobile standby does not solve the time-zero obstruction at range 1.65 without advance identity or a temporal outage allowance.

\paragraph{Separate limit at two different times.}
If the standby must cover either loss 3 at time zero or loss 1 at time 26.45, the second loss again requires reaching its singleton robot 0, now at a different pose. The exact same-position bound is half the separation between robot 4 at the first cut and robot 0 at the second. A midpoint witness reaches both components in each case and clears bodies at the two samples, giving a geometric range enclosure $[2.0538875,2.0538877]$. This does not verify that midpoint's clearance through the whole nominal path, so only the lower bound transfers to an all-time clearance contract. At the original range the two cut positions of any mobile standby would need to differ by more than $0.8077750$; no collision-free mobility plan is inferred from that bound.

\paragraph{Checks.}
The independent replay repeats \href{https://arxiv.org/html/2609.29939v1}{the published managed formation} and uses exact rational comparisons on recorded binary-point positions, strict original geometric edges and explicit squared-distance clearance tests. The logic uses the triangle inequality and the actual singleton components, without a search over fixed points or a claim about the best optimizer. There are no new vehicle, radio or clock measurements. The result separates per-fault placement with advance knowledge from a common placement and should not be used as an estimate of real communications failure. No post-fault dwell, I2 independent mine test or full R6 acceptance has been performed.
\end{document}
